% Revised Results and Discussion working draft
% Short statements are intentional. Expand only where marked TODO.

\chapter{Results and Discussion}
\label{ch:results}

This chapter answers the three research questions. RQ1 compares top-height prediction under
spatial-block cross-validation, followed by temporal and diagnostic analyses. RQ2 examines
persistent plot-level departures from a shared Chapman--Richards (CR) curve. RQ3 examines
plot-specific asymptotic deviations using global and spatially varying models. Set~3 is used for
the main RQ1 comparison as a common representative feature set across all model families,
providing a shared terrain- and wind-conditioned input set. Full feature-set results are reported
in Appendix~\ref{app:full-results}.

% Checked: Set3 first appears (as "nested_set3_gated_terrain_wind_vif") in the XGBoost
% hyperparameter search script, dated around the same time RQ1 results were being computed --
% no evidence it was fixed before test performance was inspected. Wording above uses the
% conservative "representative feature set" framing rather than a causal "because" justification.


%=======================================================================================
\section{RQ1: Prediction of top height}
\label{sec:results-rq1}

\subsection{Spatial-block performance}

Tuned XGBoost achieved the lowest mean RMSE and MAE, and the highest mean $R^2$, in both cohorts
under five-fold spatial-block cross-validation (Table~\ref{tab:results-rq1}). The DNN was the
strongest neural model in the four-survey cohort. The DNN and PINN had similar point estimates in
the six-survey cohort.

% Verified: Table~\ref{tab:results-rq1} matches the corrected/verified results table (Set~3,
% spatial\_block\_kfold, seed 42, tuned XGBoost from the 27-config search) exactly -- no
% discrepancy found against an earlier table.

\begin{table}[!htbp]
\centering
\scriptsize
\setlength{\tabcolsep}{3pt}
\renewcommand{\arraystretch}{0.95}
\begin{tabular}{lccc|ccc}
\toprule
& \multicolumn{3}{c|}{Four-survey}
& \multicolumn{3}{c}{Six-survey} \\
Model
& RMSE$\downarrow$ & MAE$\downarrow$ & $R^2\uparrow$
& RMSE$\downarrow$ & MAE$\downarrow$ & $R^2\uparrow$ \\
\midrule
Linear
& $5.168{\pm}0.361$ & $4.056{\pm}0.343$ & $0.580{\pm}0.026$
& $4.877{\pm}0.643$ & $3.619{\pm}0.426$ & $0.509{\pm}0.067$ \\
RF
& $4.794{\pm}0.187$ & $3.611{\pm}0.142$ & $0.638{\pm}0.011$
& $4.062{\pm}0.249$ & $3.010{\pm}0.185$ & $0.656{\pm}0.046$ \\
XGBoost (tuned)
& $\mathbf{4.548{\pm}0.137}$ & $\mathbf{3.431{\pm}0.123}$ & $\mathbf{0.674{\pm}0.009}$
& $\mathbf{3.656{\pm}0.232}$ & $\mathbf{2.669{\pm}0.183}$ & $\mathbf{0.722{\pm}0.031}$ \\
DNN
& $4.681{\pm}0.230$ & $3.533{\pm}0.201$ & $0.655{\pm}0.016$
& $3.903{\pm}0.316$ & $2.867{\pm}0.275$ & $0.684{\pm}0.031$ \\
PINN
& $5.206{\pm}0.399$ & $4.063{\pm}0.407$ & $0.573{\pm}0.036$
& $3.886{\pm}0.290$ & $2.819{\pm}0.196$ & $0.686{\pm}0.038$ \\
PINN-$k$
& $5.197{\pm}0.402$ & $4.046{\pm}0.421$ & $0.575{\pm}0.037$
& $3.895{\pm}0.282$ & $2.816{\pm}0.196$ & $0.684{\pm}0.038$ \\
\bottomrule
\end{tabular}
\caption{Five-fold spatial-block performance using Set~3. Values are the mean $\pm$ standard
deviation of the five fold-level metrics. Each compartment appears in one test fold. The pooled
held-out samples contain 232{,}292 plot--survey rows from 58{,}073 plots in the four-survey cohort
and 82{,}614 rows from 13{,}769 plots in the six-survey cohort.}
\label{tab:results-rq1}
\end{table}

XGBoost outperformed the DNN in both cohorts: RMSE was 0.133~m (2.8\%) lower on 4survey and
0.247~m (6.3\%) lower on 6survey. The advantage was larger over the physics-guided models in the
four-survey cohort. This is consistent with previous findings that boosted trees remain strong on
heterogeneous tabular data \citep{shwartzziv2021Tabular_C4_NN}. Possible reasons include nonlinear
thresholds, feature interactions and the limited number of spatially independent compartments.
These mechanisms were not isolated here.

Tuning effort was not equal between the two: XGBoost was selected from a 27-configuration grid
search (n\_estimators, max\_depth, learning\_rate), while the DNN used a single fixed architecture
with no formal hyperparameter search. Part of XGBoost's advantage may therefore reflect a larger
tuning budget rather than a purely architectural difference, and this should be treated as a
limitation on the RQ1 model comparison, not just on the RQ1 tuning procedure.

The four-survey cluster-bootstrap intervals did not overlap for the DNN and PINN
(DNN: $[0.624,0.690]$; PINN: $[0.545,0.607]$). This supports a meaningful performance difference.
The six-survey intervals overlapped substantially (DNN: $[0.658,0.740]$; PINN:
$[0.671,0.736]$). The six-survey results therefore do not distinguish these models reliably.

These are pooled $R^2$ 95\% intervals from 1{,}000 cluster-bootstrap resamples of whole
compartments (resampled with replacement); the interval is the 2.5th--97.5th percentile of the
resampled pooled $R^2$.


\subsection{Cohort and temporal sensitivity}

Model behaviour differed between cohorts. The six-survey cohort covers fewer compartments and a
narrower elevation range (48 compartments, elevation $102.3\pm47.9$m, range 18--351m) than the
four-survey cohort (296 compartments, elevation $177.9\pm104.6$m, range 18--561m). Its pooled CR
fit also has a lower asymptote and faster growth rate ($y_{\max}=44.46$m, $k=0.0233$ vs.
$y_{\max}=51.96$m, $k=0.0103$). These differences are consistent with sampling a smaller and more
homogeneous part of Aberfoyle. They do not isolate the cause of the model-performance difference.

Performance decreased under temporal evaluation. The reduction was larger in the six-survey
cohort. This indicates weaker transfer beyond the training period, but cohort size, coverage and
composition also differ.

\begin{table}[!htbp]
\centering
\small
\renewcommand{\arraystretch}{1.2}
\setlength{\tabcolsep}{4pt}
\begin{tabular}{llrrr}
\toprule
Cohort & Model & Spatial-block $R^2$ & Temporal $R^2$ & Absolute change \\
\midrule
4survey & DNN & 0.634 & 0.544 & $-0.090$ \\
4survey & PINN & 0.584 & 0.290 & $-0.294$ \\
4survey & PINN$_k$ & 0.588 & 0.289 & $-0.299$ \\
6survey & DNN & 0.722 & 0.317 & $-0.405$ \\
6survey & PINN & 0.731 & 0.161 & $-0.570$ \\
6survey & PINN$_k$ & 0.730 & 0.190 & $-0.540$ \\
\bottomrule
\end{tabular}
\caption{Spatial-block vs.\ temporal-holdout $R^2$, by model and cohort. Training uses survey data
up to 2012, spatial-block validation uses 2021 data, temporal holdout tests on 2023 data. Only
$R^2$ was computed for this check; RMSE and MAE are not available from this source
\fc{(re-run if RMSE/MAE are needed for the final draft)}.}
\label{tab:results-rq1-temporal}
\end{table}

The PINN and PINN$_k$ lost far more accuracy under temporal evaluation than the DNN, on both
cohorts. If the physics constraint improved extrapolation beyond the training period, the opposite
pattern would be expected: a smaller temporal drop for the physics-guided models, not a larger one.
This is one further respect in which CR guidance did not deliver its intended benefit in this
setting, alongside its lack of a spatial-block accuracy advantage (Table~\ref{tab:results-rq1}).
It does not by itself explain why the constraint fails to generalise better in time; cohort
composition and the widening gap between training years and the 2023 test survey are also
plausible contributors and were not separated here.


\subsection{Effect of the evaluation split}

The DNN achieved substantially higher $R^2$ under the random plot-level split than under the
spatial split (4survey: $0.831$ vs.\ $0.634$, a gap of $0.197$; 6survey: $0.844$ vs.\ $0.722$, a
gap of $0.122$). Nearby training and test plots can share similar stand and environmental
conditions, so the random split gives an optimistic estimate of spatial transfer. The PINN's
random-split and spatial-split $R^2$ were nearly identical (4survey: $0.590$ vs.\ $0.584$;
6survey: $0.728$ vs.\ $0.731$), a near-zero gap rather than a smaller one. This difference in gap
size is associative, not proven causal: the PINN's physics constraint, its generally weaker fit,
and other architectural differences between the two models were not tested separately, so no
single factor can be isolated as the cause. These results support the use of compartment-level
cross-validation as the primary evaluation.


\subsection{Physics-weight ablation}

Increasing the physics and trajectory weights reduced predictive performance. The decrease was
largest in the four-survey cohort. The effect was smaller in the six-survey cohort.

Physical guidance also changed the learned CR parameters. However, the $\lambda=0$ parameter
sub-networks received no optimisation signal and collapsed to fold-level constants. Their
$y_{\max}$--$k$ correlation is therefore not a valid plot-level identifiability control. The
results show that physical guidance affected parameter estimates, but do not establish that it
resolved identifiability relative to an otherwise equivalent unconstrained parameter model.

% TODO: Add fold-level parameter summaries for lambda > 0 if they are reliable.
% TODO: If retained, report the classical CR y_max--k correlation only as contextual evidence.
% TODO: Move the invalid lambda=0 plot-level CI and bootstrap calculation out of the main chapter.

% FIGURE 1 SPECIFICATION
% Chart: two-panel paired fold plot.
% Panel A: four-survey; Panel B: six-survey.
% x-axis: model; y-axis: fold R2. Join values from the same fold with thin grey lines.
% Highlight XGBoost, DNN and PINN; mute other baselines.
% Why: shows whether the headline ranking holds in each spatial fold. Adds information not visible
% in the mean+-SD table. Avoid another bar chart of the same averages.
\begin{figure}[!htbp]
\centering
\fbox{\parbox{0.92\textwidth}{\centering\vspace{3.0cm}
TODO FIGURE: paired fold-level $R^2$ comparison for XGBoost, DNN and PINN; one panel per cohort
\vspace{3.0cm}}}
\caption{Fold-level spatial performance of the main competing models. Thin lines connect models
evaluated on the same held-out compartment fold. This shows the consistency of the model ranking,
rather than repeating the table means.}
\label{fig:rq1-fold-performance}
\end{figure}

% FIGURE 2 SPECIFICATION
% Chart: line plot of lambda against held-out R2.
% Separate lines for PINN and PINN-k; panels for cohorts.
% Optional lower panel: a valid parameter-stability measure for lambda > 0 only.
% Why: directly shows the accuracy--constraint trade-off and whether it is monotonic.
% Do not plot the invalid lambda=0 parameter correlation as a plot-level estimate.

\paragraph*{Answer to RQ1.}
Tuned XGBoost gave the strongest spatially held-out predictions in both cohorts. The DNN remained
competitive, but CR guidance did not improve predictive accuracy. Random plot splitting gave an
optimistic DNN estimate. The current $\lambda=0$ implementation does not support a controlled
claim about parameter identifiability.


%=======================================================================================
\subsection{Diagnostic: where environmental models improve on the shared curve}
\label{sec:results-rq1-cr-reduction}

Environmental models did not improve every observation uniformly. Error reductions were
concentrated among observations with the largest CR baseline errors. Performance often worsened
in the best-fitting quartile.

\begin{table}[!htbp]
\centering
\scriptsize
\renewcommand{\arraystretch}{1.1}
\begin{tabular}{llrrrr}
\toprule
Model & Cohort & CR MAE (m) & Reduction (m) & Reduction (\%) & Rows improved (\%) \\
\midrule
DNN, environment & Four-survey & 5.028 & $1.501\pm0.023$ & 29.9 & $64.8\pm4.1$ \\
DNN, environment & Six-survey & 3.047 & $0.174\pm0.036$ & 5.7 & $52.4\pm4.8$ \\
XGBoost, environment & Four-survey & 5.028 & $1.597\pm0.330$ & 31.8 & $67.2\pm3.8$ \\
XGBoost, environment & Six-survey & 3.047 & $0.361\pm0.183$ & 11.8 & $56.0\pm4.6$ \\
XGBoost, no environment & Four-survey & 5.030 & $1.222\pm0.232$ & 24.3 & $64.5\pm3.2$ \\
XGBoost, no environment & Six-survey & 3.047 & $0.240\pm0.144$ & 7.9 & $54.8\pm3.0$ \\
\bottomrule
\end{tabular}
\caption{Reduction in absolute error relative to the fold-matched CR baseline. Positive values
indicate lower error than CR. The DNN uncertainty summarises five seeds; XGBoost uses one fixed
configuration. These uncertainty terms are therefore not directly equivalent.}
\label{tab:results-cr-reduction}
\end{table}

The strongest reductions occurred in Q4, the quartile with the largest CR errors: the model cuts
Q4 error by 5.08~m on 4survey and 1.81~m on 6survey. Q1 often became less accurate, worsening by
$2.6\times$ on 4survey and $3.2\times$ on 6survey. However, the quartiles are defined using CR
error. Q4 has more error available to reduce, while Q1 has little scope for improvement. This
analysis shows where replacement models differ from CR. It does not alone prove selective
environmental correction.

The no-environment XGBoost control retained part of the improvement. Environmental inputs added a
further reduction of 0.375~m in the four-survey cohort and 0.121~m in the six-survey cohort.

The environmental predictors were shuffled as one block, jointly rather than column-by-column, so
inter-predictor correlations were preserved; the model was then retrained on the shuffled data
(a training-time, not prediction-time, permutation), with five repeats, one per spatial-block
fold. Shuffling the environmental predictors left most Q1 degradation intact but reduced the Q4
gain, reported as the mean reduction across folds $\pm$ its standard deviation. This is consistent
with two components: a general cost from replacing an already accurate CR prediction, and an
informative contribution from unshuffled environmental variables among poorly fitted observations.
This is the evidence that separates the quartile pattern above from a purely mechanical artefact of
the quartile definition: if Q4's gain were only "more room to improve," it should survive shuffling
as well as Q1's degradation does, but it does not.

No reliable spatial clustering of error reduction was detected. The six-survey semivariogram
remained flat within the tested range. In the four-survey cohort, Moran's $I$ approached zero and
changed sign as the distance window increased. These results do not support a stable geographical
``help map''.

% Move exact windows, sign changes and the 47-compartment subsampling check to the appendix.

% FIGURE 3 SPECIFICATION
% Chart: dot-and-line interaction plot, not grouped bars.
% x-axis: CR-error quartile Q1--Q4; y-axis: mean absolute-error reduction in metres.
% Panels: four-survey and six-survey.
% Lines: environment model, no-environment control, shuffled-environment control.
% Add a horizontal zero line. Positive = improvement; negative = harm.
% Why: shows the central trade-off, model controls and cohort replication in one figure. The table
% only gives the aggregate and misses this mechanism.
\begin{figure}[!htbp]
\centering
\fbox{\parbox{0.92\textwidth}{\centering\vspace{3.0cm}
TODO FIGURE: CR-error quartile against mean error reduction; controls and cohorts shown separately
\vspace{3.0cm}}}
\caption{Error reduction relative to CR across quartiles of CR baseline error. The zero line
separates improvement from degradation. This figure shows whether gains are concentrated among
observations that depart most from the shared curve.}
\label{fig:cr-error-quartiles}
\end{figure}


%=======================================================================================
\section{RQ2: Persistent departure from the shared curve}
\label{sec:results-rq2}

RQ2 models each plot's mean residual across its four-survey history. Canopy cover and thinning
variables had the largest attribution values across the tested feature sets and models.

\begin{table}[!htbp]
\centering
\scriptsize
\renewcommand{\arraystretch}{1.1}
\begin{tabular}{lcccc}
\toprule
Set & Elastic Net $R^2$ [95\% CI] & XGBoost $R^2$ [95\% CI]
& EN Moran's $I$ & XGB Moran's $I$ \\
\midrule
Set1 & 0.220 [0.180, 0.252] & 0.222 [0.183, 0.252] & 0.145 & 0.146 \\
Set2 & 0.359 [0.310, 0.396] & 0.416 [0.355, 0.461] & 0.105 & 0.032 \\
Set3 & 0.325 [0.272, 0.365] & 0.375 [0.321, 0.414] & 0.110 & 0.077 \\
Set4 & 0.358 [0.305, 0.398] & 0.395 [0.338, 0.438] & 0.103 & 0.054 \\
\bottomrule
\end{tabular}
\caption{Held-out attribution of each plot's mean CR residual in the four-survey cohort. Moran's
$I$ is evaluated using each model's semivariogram-resolved distance band. Full ranges, permutation
tests and mixed-effects results are reported in Appendix~\ref{app:full-results}.}
\label{tab:results-rq2}
\end{table}

Environmental feature sets improved held-out attribution over the structural baseline: every
feature set's interval clears the baseline-only Set~1 without overlap. Set~2 had the highest
XGBoost point estimate. Its interval overlapped Set~4, so the results do not establish a
difference between these feature sets. Terrain-only Set~3, which excludes soil and climate
predictors, consistently had the lowest point estimates and interval bounds of the three -- soil
and climate variables therefore add attribution signal beyond terrain and wind exposure alone, and
a terrain-only feature set is not a sufficient substitute.

Predictors were standardised (mean 0, SD 1) before fitting Elastic Net and the mixed-effects
model, so their coefficients are metres of mean CR residual per one-SD change in each predictor.
XGBoost, being tree-based, does not require standardisation; its SHAP values are reported on the
original response scale, so magnitudes are not directly comparable to the Elastic Net
coefficients above.

Canopy cover had the largest Elastic Net coefficient and XGBoost SHAP magnitude. Removing it from
Set~4 reduced held-out $R^2$ in both models \sr{(exact removal-$R^2$ values not added: no ablation
script matching Set~4's folds and final model version was found this session -- verify before
citing a number)}. Canopy cover was also moderately correlated with the height response. Its
importance may reflect stand structure, prior growth and the shared LiDAR source. It is therefore
treated as a structural control, not an independent environmental driver.

Slope was the only abiotic predictor with a stable direction and substantial magnitude across the
principal global models. TOPEX was positive in the mixed-effects model but weak and unstable in
Elastic Net. Because TOPEX varies mainly between compartments, its coefficient may be sensitive to
unmeasured compartment differences and fold composition. This is a plausible explanation, not an
isolated causal mechanism.

XGBoost left less residual spatial autocorrelation than Elastic Net as environmental features were
added, though both start from a similar low-to-moderate baseline. Moran's $I$ falls from 0.146
(Set~1) to 0.054 (Set~4) for XGBoost, a 63\% relative drop, against only 0.145 to 0.103 for
Elastic Net, a 29\% relative drop. \ai{Plausible explanation: Elastic Net is linear and additive,
so it can only remove the linear share of each predictor's spatial signal; any nonlinear or
interaction effect (e.g.\ a terrain/wind threshold) stays in its residual, and terrain-derived
predictors are themselves spatially smooth, so that leftover signal still clusters by location.
XGBoost's tree structure fits such effects directly, removing more of the spatially structured
signal. This is the same model-flexibility gap the RQ1 XGBoost-vs-DNN result raises -- stated here
as a hypothesis consistent with the pattern, not directly tested.} \lim{Adding interaction or
polynomial terms to Elastic Net (or fitting a GAM with the same predictors) would directly test
this, rather than leaving it untested.}

% Checked: the second-stage model uses statsmodels MixedLM (models/spatial_attribution/nlme.py),
% a LINEAR mixed model fit to CR residuals (fixed terrain/wind effects, random compartment
% intercept) -- not a genuinely nonlinear mixed-effects fit. "mixed-effects model" (used
% throughout this section) is the correct label; "NLME" should not be used for it.
% TODO: Define the two-stage mixed-effects output in the appendix (fixed/random effect structure,
% what "mixed-effects model" refers to here) -- exact appendix location not established this
% session, content above is ready to use.

% FIGURE 4 SPECIFICATION
% Composite with two connected panels.
% Panel A: coefficient/SHAP rank plot for only the 6--8 most important features. Show direction for
% Elastic Net and mixed effects; magnitude for SHAP. Mark predictors with stable sign across folds.
% Panel B: residual Moran's I by Set for EN and XGBoost, with connected points.
% Why: Panel A answers what explains departure. Panel B shows what remains spatially unexplained.
% A residual map may be a small inset only if it reveals a pattern not conveyed by Moran's I.

\paragraph*{Answer to RQ2.}
Stand structure and management explained more persistent CR departure than any single abiotic
predictor. Slope was the most stable abiotic association. Environmental features improved held-out
attribution, but residual spatial autocorrelation remained.


%=======================================================================================
\section{RQ3: Plot-specific asymptotic deviation}
\label{sec:results-rq3}

RQ3 examines the difference between each plot's fitted asymptote and its yield-class expectation.
It compares global attribution models with GNNWR, which allows relationships to vary by location.

\begin{table}[!htbp]
\centering
\scriptsize
\renewcommand{\arraystretch}{1.1}
\begin{tabular}{lcc|cc|cc}
\toprule
& \multicolumn{2}{c|}{Elastic Net $R^2$}
& \multicolumn{2}{c|}{XGBoost $R^2$}
& \multicolumn{2}{c}{GNNWR $R^2$} \\
Set & Four-survey & Six-survey & Four-survey & Six-survey & Four-survey & Six-survey \\
\midrule
Set2 & 0.286 [0.241, 0.325] & 0.057 & 0.266 [0.214, 0.308] & 0.031
& 0.320 [0.270, 0.369] & $0.014\pm0.053$ \\
Set3 & 0.274 [0.223, 0.316] & 0.031 & 0.281 [0.235, 0.320] & 0.086
& 0.319 [0.263, 0.365] & $0.054\pm0.055$ \\
Set4 & 0.240 [0.192, 0.286] & 0.056 & 0.250 [0.201, 0.295] & 0.015
& 0.294 [0.236, 0.347] & $0.002\pm0.068$ \\
\bottomrule
\end{tabular}
\caption{Held-out $R^2$ for plot-specific asymptotic deviation. Four-survey brackets are 95\%
cluster-bootstrap intervals. Six-survey GNNWR values are the mean $\pm$ standard deviation across
three split seeds. Other six-survey entries are single estimates. These uncertainty summaries are
not directly equivalent.}
\label{tab:results-rq3}
\end{table}

GNNWR had the highest four-survey point estimate in all three feature sets. Its confidence
intervals overlapped those of the global models. The results therefore suggest, but do not
establish, a predictive advantage. GNNWR also left lower residual Moran's $I$ in every
four-survey feature set, indicating less residual spatial clustering.

All models performed poorly on the six-survey cohort. No stable residual spatial structure was
detected within the tested distance windows. This prevents Moran's $I$ from explaining the
near-zero predictive performance. Matching the compartment count through four-survey subsampling
did not reproduce the same behaviour. Cohort composition may therefore matter more than
compartment count alone.

% TODO: move the block below into Appendix~\ref{app:full-results} (exact appendix file/location
% not established in this session -- content is ready, placement is not).
% Appendix content -- exact semivariogram ranges, p-values, 10km windows, subsample values:
%   RQ2 (mean CR residual, 4survey): ranges 1665-1739m (EN), 1385-1675m (XGB), all p=0.001.
%   RQ3 (asymptote deviation): 4survey ranges 808-1252m, all p=0.001. 6survey: Set3 unresolved
%     (flat semivariogram) for EN/XGB/GNNWR; Set2/Set4 XGB+GNNWR only resolve after widening to
%     10000m, landing at I~-0.0002, p=0.05-0.31 (not significant); 6survey EN unresolved in all
%     3 sets.
%   47-compartment subsample check (4survey, matching 6survey's compartment count): 3 random
%     seeds, all resolve cleanly, I=0.18-0.36, ranges 1045-1659m -- rules out compartment count
%     alone as the explanation for 6survey's unresolvability.

Canopy cover ranked highly in the four-survey global and spatial models. Other feature rankings
varied between methods. These differences should be interpreted cautiously because each method
represents spatial structure differently. Six-survey feature importance is not interpreted here:
held-out $R^2$ for that cohort is near zero (Table~\ref{tab:results-rq3}, $0.014$--$0.086$), so any
ranking would reflect noise rather than a genuine predictor-outcome relationship.


\subsection{Extreme asymptotic deviations}

Tukey fences flagged 266 four-survey plots (0.47\%) and 709 six-survey plots (5.32\%). Their raw
height trajectories were often plausible, but some fitted asymptotes reached implausible values.
Fixing the growth-rate and shape parameters can force the fitted asymptote upward when a plot grows
more slowly than its assigned yield-class curve. Standard raw-height filters do not detect this
target-construction instability.

Removing flagged plots did not improve predictive performance. Four-survey $R^2$ changed by only
$-0.002$. Six-survey performance decreased. Exclusion is therefore not supported solely by model
accuracy, but the fitted asymptotes should not be interpreted as literal biological ceilings.

None of the 266 flagged four-survey plots met the clearfelling or measurement-anomaly rules
\fc{(exact rule definitions still need writing up in the methodology/appendix, not yet
documented)}. This rules out active disturbance as the source of these extreme values and
favours the target-construction explanation above: the flagged asymptotes are an artefact of how
the target is built, not evidence of real biological ceilings or unrecorded felling events.
Flagged plots sit closer to compartment boundaries than the population (median 9.4~m vs.\ 28.0~m)
and were disproportionately poorly fitted in 2008 (75.6\% of flagged plots fit worst in that year,
vs.\ 47.8\% baseline), pointing toward boundary and yield-class assignment error specifically,
rather than a single dominant cause. Possible explanations include clipped plot geometry, boundary
changes, mixed stand conditions and imperfect yield-class assignment. These mechanisms were not
isolated.

% FIGURE 5 SPECIFICATION
% Three-panel explanatory figure.
% A: map of asymptotic deviation, with flagged plots outlined rather than dominating the colour
% scale. Use a robust symmetric colour scale based on central quantiles.
% B: three example height--age trajectories: stable fit, moderate deviation, unstable high
% asymptote. Plot observed heights, fixed-shape fitted curve and yield-class reference.
% C: compact distribution of deviation by distance-to-boundary group or flagged status.
% Why: explains what the target means, why extremes occur, and where they occur. This is more
% informative than another model-score chart.
\begin{figure}[!htbp]
\centering
\fbox{\parbox{0.92\textwidth}{\centering\vspace{3.0cm}
TODO FIGURE: deviation map, three trajectory archetypes, and boundary comparison
\vspace{3.0cm}}}
\caption{Spatial distribution and construction of plot-specific asymptotic deviations. Example
trajectories distinguish plausible departures from high asymptotes produced by the fixed-shape
fit. Flagged plots remain in the main analysis.}
\label{fig:rq3-deviation-examples}
\end{figure}

% OPTIONAL FIGURE 6 SPECIFICATION
% Chart: slope chart or dot plot of residual Moran's I for EN, XGBoost and GNNWR across Sets 2--4.
% Four-survey only. Pair models within each set. Include semivariogram range in caption/table, not
% encoded in the same chart.
% Why: supports the spatial-modelling claim without repeating R2 values.

\paragraph*{Answer to RQ3.}
GNNWR produced the highest four-survey point estimates and the lowest residual spatial
autocorrelation, but overlapping intervals do not establish a clear predictive advantage. None of
the models generalised well on the six-survey cohort. Extreme target values were often consistent
with instability from the fixed-shape asymptote construction and should not be treated as literal
ecological ceilings.


%=======================================================================================
\section{Overall discussion}
\label{sec:results-overall-discussion}

Three findings connect the research questions. First, tuned XGBoost provided the strongest
top-height predictions. Second, CR guidance reduced predictive accuracy under the tested loss
weights. Third, stand structure and management explained more curve departure than any single
abiotic predictor, although spatial structure remained.

The weak PINN result does not show that CR theory is unsuitable for forest growth. It shows that
the tested derivative and trajectory penalties did not improve held-out prediction for these
data. A pooled CR derivative may constrain local trajectories that differ through management,
disturbance or site conditions. The DNN may also capture most of the predictable age--environment
relationship without the additional penalty.

The attribution results describe association, not causation. The data cover one forest. Several
environmental variables are spatially coarse or model-derived. Management records are incomplete.
Canopy cover shares its LiDAR source with the response. These limits restrict ecological
interpretation and transfer beyond Aberfoyle.

% TODO: Expand by one paragraph only:
% - practical meaning for forest monitoring;
% - why spatial validation changes the conclusion;
% - what a future external-site test should evaluate.
% Keep detailed future experiments for the Conclusion/Future Work chapter.

